% p081_1; source PDF page 81; manually verified editable table.
\begin{tikzpicture}[x=1cm,y=-1cm,inner sep=0pt,outer sep=0pt]
\path[use as bounding box] (0,0) rectangle (16.000,7.963);
\draw[line width=.4pt] (0.000,0.000) rectangle (3.400,1.048);
\node[align=center,text width=3.140cm,font={\linespread{1}\fontsize{10}{13.50}\selectfont\bfseries}] at (1.700,0.524) {关于\CJKunderwave{犯罪既遂的条件}};
\draw[line width=.4pt] (3.400,0.000) rectangle (7.000,1.048);
\node[align=center,text width=3.340cm,font={\linespread{1}\fontsize{10}{13.50}\selectfont\bfseries}] at (5.200,0.524) {因果关系};
\draw[line width=.4pt] (7.000,0.000) rectangle (16.000,1.048);
\node[align=center,text width=8.740cm,font={\linespread{1}\fontsize{10}{13.50}\selectfont\bfseries}] at (11.500,0.524) {范例};
\draw[line width=.4pt] (0.000,1.048) rectangle (3.400,4.958);
\node[align=center,text width=3.140cm,font={\linespread{1}\fontsize{10}{13.50}\selectfont}] at (1.700,3.003) {客观要件};
\draw[line width=.4pt] (3.400,1.048) rectangle (7.000,3.553);
\node[align=left,text width=3.340cm,font={\linespread{1}\fontsize{10}{13.50}\selectfont}] at (5.200,2.300) {没有因果关系，不构成故意犯罪既遂};
\draw[line width=.4pt] (7.000,1.048) rectangle (16.000,3.553);
\node[align=left,text width=8.740cm,font={\linespread{1}\fontsize{10}{13.50}\selectfont}] at (11.500,2.300) {（人参果案）孙悟空想毒杀猪八戒，给其有毒的人参果。猪八戒吃时卡在喉咙，噎死。毒药没有进入胃里起作用，毒杀行为与死亡没有因果关系，不构成故意杀人罪既遂，而是未遂①（2023 年试题）};
\draw[line width=.4pt] (3.400,3.553) rectangle (7.000,4.958);
\node[align=left,text width=3.340cm,font={\linespread{1}\fontsize{10}{13.50}\selectfont}] at (5.200,4.255) {具有因果关系，进入主观要件$\to$};
\draw[line width=.4pt] (0.000,4.958) rectangle (3.400,7.963);
\node[align=center,text width=3.140cm,font={\linespread{1}\fontsize{10}{13.50}\selectfont}] at (1.700,6.460) {主观要件};
\draw[line width=.4pt] (3.400,4.958) rectangle (7.000,7.963);
\node[align=left,text width=3.340cm,font={\linespread{1}\fontsize{10}{13.50}\selectfont}] at (5.200,6.460) {以为是 A 死法，实际是 B 死法。\\{}[结论]构成故意犯罪既遂};
\draw[line width=.4pt] (7.000,3.553) rectangle (16.000,7.963);
\node[align=left,text width=8.740cm,font={\linespread{1}\fontsize{10}{13.50}\selectfont}] at (11.500,5.758) {（推下井案）甲将乙推入井里，欲淹死乙，井里没水，摔死了乙。客观上，甲的行为与乙的死亡存在因果关系。主观上，甲关于具体因果流程的认识错误不重要。甲构成故意杀人罪既遂。（2011 年第 53 题）\\{}[总结]人参果案与推下井案的相同点是，都以为是 A 死法导致死亡，实际是 B 死法导致死亡。区别在于，前者在客观要件中没有因果关系；后者在客观要件中有因果关系};
\end{tikzpicture}
